\section{Background and Added Value}

My MSc thesis at Politecnico di Milano, where I built, calibrated, and ran sensitivity and uncertainty analysis on the \textbf{FEST} process-based hydrological model over the Po basin, gave me direct experience of the central methodological tension that this project addresses: constraining subsurface parameters (hydraulic conductivity, deep percolation factors, \textbf{SCS} curve number) from surface observations alone.

My current operational role at Progesi S.p.A., where I run multi-hazard monitoring with \textbf{ICON} and \textbf{ECMWF IFS} forcing, satellite-derived diagnostics, and Python, R, and MATLAB workflows for Regione Lombardia, translates directly to the operational demands of a national-scale drought early warning system.

The combination of process-based modelling, applied data science, and recent exposure to operational hydrometeorology positions me to start contributing to \textsc{WP1} and \textsc{WP2} immediately, while developing the hybrid machine learning competence required for \textsc{WP3} through structured upskilling in \textbf{NeuralHydrology}, PyTorch, and differentiable hydrology.
